% !TEX program = xelatex
% 编译方式：把本文件夹整包复制到纯 ASCII 路径后，执行两遍 xelatex（详见 docs/_manual_common/README.md）。
% 首版骨架：内容取自 src/short_circuit/ 与 include/hacdcpf/analysis/ 源码
%           （short_circuit.hpp / dc_short_circuit.hpp），公式与参数以代码为准。
\documentclass[UTF8,fontset=windows,a4paper,11pt]{ctexart}
\usepackage{hysim_manual}

\title{\textbf{交直流混合高弹性能源电力系统仿真平台}\\[0.5em]
       {\Large 短路计算技术手册（IEC 60909 与直流故障水平）}\\[0.3em]
       {\large 数学模型、接口参数与验证校核（首版骨架）}}
\author{HySim-XJTU-HRPES（西安交通大学 HRPES 团队）}
\date{\today}

\begin{document}
\maketitle
\begin{abstract}
\noindent 本手册依据 HySim-XJTU-HRPES 平台短路模块
（\srcpath{src/short\_circuit/}、公共头 \srcpath{include/hacdcpf/analysis/short\_circuit.hpp} 与
\srcpath{dc\_short\_circuit.hpp}）整理，覆盖两条路径：交流侧基于 IEC 60909 等效电压源的
序网 Z-bus 方法（概览与详细），给出 $I_k''$、峰值 $i_p$、开断 $I_b$、稳态 $I_k$、热等效
$I_{th}$ 及换流器贡献；直流侧为电阻性刚性源螺栓短路的保守上界估计。全部结论以源码与
\srcpath{docs/short\_circuit\_rich\_acdc\_derivation.md} 为准：交流分析在
\srcpath{project\_to\_canonical\_models(sys)} 之后进行；本模块不是保护配合仿真器。
文档版式与《碳流追踪与碳排放核算技术手册》一致。
\end{abstract}

\tableofcontents
\newpage

\section{阅读指南与约定}
\subsection{数据来源}
\begin{itemize}
  \item \srcpath{include/hacdcpf/analysis/short\_circuit.hpp}、
        \srcpath{src/short\_circuit/short\_circuit.cpp} —— 交流短路接口与实现；
  \item \srcpath{include/hacdcpf/analysis/dc\_short\_circuit.hpp}、
        \srcpath{src/short\_circuit/dc\_short\_circuit.cpp} —— 直流故障水平；
  \item \srcpath{docs/short\_circuit\_rich\_acdc\_derivation.md} —— 交直流短路推导契约；
  \item \srcpath{tests/test\_short\_circuit\_iec60909\_4.cpp} 等 —— 验证依据。
\end{itemize}

\subsection{单位制与索引空间}
电流以标幺值（pu）与 kA 并列输出（\fld{i\_fault\_pu}/\fld{ikpp\_ka}、\fld{i\_fault\_ka}），
短路容量 $S_k$ 单位 MVA，阻抗/电阻标幺化，时间 \fld{ith\_duration\_s}/\fld{breaking\_time\_s}
单位 s。交流结果索引空间与规范化后的母线一致（分析前先做
\srcpath{project\_to\_canonical\_models}）；直流结果以 \texttt{DC\_V} 源母线为戴维南参考。

\subsection{诚实口径}
概览交流路径（\srcpath{compute\_short\_circuit}）对不平衡故障采用简化的 $Z_1=Z_2=Z_0$ 假设
（负序取正序、零序未建模）；详细路径构造真实正/负/零序矩阵，但三绕组零序默认以正序代入
（未建向量组）。直流路径\textbf{不}在计算内部建模断路器开合状态、换流器限流、电容放电暂态
与保护动作，仅给出受线路电阻限制的保守上界，用于断路器/电缆选型筛查。相关限制在 \S 5.2
如实列出。

\section{功能定位与架构}
\subsection{公共入口族}
\begin{center}\footnotesize
\begin{tabular}{@{}L{6.8cm}L{3.0cm}L{5.2cm}@{}}
\toprule
入口 & 定义位置 & 说明\\
\midrule
\srcpath{compute\_short\_circuit(sys, opt)} & short\_circuit.hpp & 全母线 Z-bus 概览短路\\
\srcpath{compute\_fault\_at\_bus(sys, bus\_id, opt)} & 同上 & 单母线故障查询\\
\srcpath{run\_short\_circuit\_detailed(sys, bus, opt)} & 同上 & 详细 IEC 60909（$I_k''$/$i_p$/$I_b$/$I_k$/$I_{th}$）\\
\srcpath{run\_short\_circuit\_detailed\_batch(...)} & 同上 & 多母线详细计算\\
\srcpath{dc\_bus\_fault\_level(sys, dc\_bus\_id, opt)} & dc\_short\_circuit.hpp & 直流螺栓短路电阻性估计\\
\srcpath{get\_voltage\_factor\_sc / calculate\_kappa\_basic\_sc} & short\_circuit.hpp & IEC 电压系数与 $\kappa$ 辅助函数\\
\bottomrule
\end{tabular}
\end{center}

\subsection{关键数据结构}
\compmeta{SCDetailedResult}{include/hacdcpf/analysis/short\_circuit.hpp}
主要字段：\fld{fault\_bus\_id}、\fld{solved}、\fld{bus\_results}
（\texttt{SCDetailedBusResult}：\fld{ikss\_ka}、\fld{ip\_ka}、\fld{ib\_ka}、\fld{ik\_ka}、
\fld{ith\_ka}、\fld{v\_remaining\_pu} 及各源贡献）、\fld{branch\_results}、
\fld{converter\_contributions}。概览结果 \texttt{SCResult}/\texttt{BusFaultResult} 含
\fld{z\_thevenin}、\fld{i\_fault\_pu}、\fld{sk\_mva}、\fld{ikpp\_ka}。

\compmeta{DCFaultResult}{include/hacdcpf/analysis/dc\_short\_circuit.hpp}
主要字段：\fld{solved}、\fld{v\_prefault\_pu}、\fld{r\_thevenin\_pu}、\fld{i\_fault\_pu}、
\fld{i\_fault\_ka}、\fld{breaker\_duties}（\texttt{DCBreakerDutyResult}）。

\section{核心数学模型}
\subsection{交流 IEC 60909 序网 Z-bus}
以等效电压源法（电压校正系数 $c$）在序网上做 Z-bus 计算，复导纳矩阵经 Eigen 复数
\texttt{SparseLU} 分解。母线 $k$ 的初始对称短路电流为
\begin{equation}
I_k'' = \frac{c\,U_n}{\sqrt{3}\,\lvert Z_{kk}\rvert},
\end{equation}
峰值电流由冲击系数 $\kappa$ 给出，$\kappa$ 随 $R/X$ 衰减
\begin{equation}
i_p = \kappa\,\sqrt{2}\,I_k'', \qquad \kappa = 1.02 + 0.98\,e^{-3R/X},
\end{equation}
热等效电流由因子 $m,n$ 合成
\begin{equation}
I_{th} = I_k''\sqrt{m+n}.
\end{equation}
换流器短路模型：并网跟随型取受限电流源 $I_{sc}=k\,I_{\text{rated}}$，构网型取阻抗后电压源。

\subsection{直流螺栓短路（电阻性上界）}
直流侧以电阻性刚性源近似，母线 $k$ 螺栓短路电流为
\begin{equation}
I_f = \frac{V_k}{R_{kk} + R_f},
\end{equation}
其中 $R_{kk}$ 为到 \texttt{DC\_V} 源母线的戴维南电阻、$R_f$（\fld{fault\_resistance\_pu}）为
故障电阻。该式给出受线路电阻限制的保守上界，用于断路器分断/电缆热定值筛查。

\section{接口与参数}
\begin{paramtable}
\fld{c\_factor} & $c$ & — & $1.0\sim1.1$ & $1.1$ & IEC 电压校正系数\\
\fld{default\_xdpp} & $X_d''$ & pu & — & $0.2$ & 默认次暂态电抗\\
\fld{compute\_all\_buses} & — & 布尔 & — & \texttt{true} & 是否计算全部母线\\
\fld{inverse\_rhs\_batch\_size} & — & — & 整数 & $32$ & 逆求解右端批大小\\
\fld{ith\_duration\_s} & $T_k$ & s & — & $1.0$ & 热等效电流持续时间\\
\fld{breaking\_time\_s} & $t_{\min}$ & s & — & $0.10$ & 断路器最小分断时间\\
\fld{fault\_resistance\_pu} & $R_f$ & pu & $\ge 0$ & $0$ & 直流故障电阻\\
\fld{source\_voltage\_pu} & — & pu & $\ge 0$ & $0$ & 直流源电压（$0$=取预故障）\\
\fld{min\_closed\_breaker\_resistance\_pu} & — & pu & 小正数 & $10^{-6}$ & 闭合直流断路器最小电阻\\
\end{paramtable}
详细路径 \srcpath{run\_short\_circuit\_detailed} 通过 \texttt{SCDetailedOptions} 选择
\fld{calc\_type}（\texttt{Max}/\texttt{Min}）、\fld{kappa\_method}（A/B/C）、\fld{topology}
（\texttt{Radial}/\texttt{Meshed}）与 \fld{apply\_iec\_transformer\_correction}；直流侧
\fld{consider\_dc\_breakers}/\fld{dc\_breakers\_control\_branches} 控制断路器支路处理。

\section{验证与边界局限}
\subsection{验证与校核}
注册测试：\srcpath{test\_short\_circuit\_crossval}（Z-bus 戴维南交叉校核与 IEC 检查，
兼测 \srcpath{run\_short\_circuit\_detailed}）、\srcpath{test\_dc\_short\_circuit}、
\srcpath{test\_short\_circuit\_iec60909\_4}（IEC TR 60909-4:2021 §6.2 参考算例），并被
\srcpath{test\_io\_etap} 引用。

\subsection{边界与局限（如实标注）}
\begin{itemize}
  \item 概览交流路径对不平衡故障用简化 $Z_1=Z_2=Z_0$（负序取正序、零序未建模）；详细路径
        建真实序网，但三绕组零序默认以正序代入（未建向量组）；
  \item 直流路径不在计算内部建模断路器开合状态、换流器限流、电容放电暂态与保护动作，
        仅为受线路电阻限制的保守上界，用于断路器/电缆选型筛查；
  \item 非保护配合仿真器：无动作时间曲线、电弧模型或跳闸后拓扑更新。
\end{itemize}
\end{document}
